\documentclass[tikz,border=8pt]{standalone}
\usepackage{ctex}
\usepackage{amsmath}
\usetikzlibrary{calc, arrows.meta}

\begin{document}
\begin{tikzpicture}[scale=0.82, thick]

% 4 种方位抛物线对比：2×2 布局
% 左列：开口右（y^2=2px）和开口左（y^2=-2px）
% 右列：开口上（x^2=2py）和开口下（x^2=-2py）
% 每个子图中心坐标：(0,0), (5.5,0), (0,-5.5), (5.5,-5.5)

%% --- 子图1：开口右 y^2 = 2px (p=2, y^2=4x) ---
\begin{scope}[shift={(0,0)}]
  \draw[->, thin, gray] (-1.5,0) -- (3.2,0) node[right, font=\tiny] {$x$};
  \draw[->, thin, gray] (0,-2.3) -- (0,2.3) node[above, font=\tiny] {$y$};
  \node[font=\tiny, below right] at (0,0) {$O$};
  % 抛物线 y^2=4x
  \draw[black, thick, domain=-2.0:2.0, smooth, samples=60]
    plot ({(\x)^2/4}, \x);
  % 焦点
  \fill[black] (1,0) circle [radius=2pt];
  \node[below right, font=\tiny] at (1,0) {$F$};
  % 准线
  \draw[red, thin] (-1,-2.1) -- (-1,2.1);
  % 标注
  \node[font=\small, align=center] at (1.0,-2.8) {$y^2=2px$\\开口向右};
\end{scope}

%% --- 子图2：开口左 y^2 = -2px (p=2, y^2=-4x) ---
\begin{scope}[shift={(5.5,0)}]
  \draw[->, thin, gray] (-3.2,0) -- (1.5,0) node[right, font=\tiny] {$x$};
  \draw[->, thin, gray] (0,-2.3) -- (0,2.3) node[above, font=\tiny] {$y$};
  \node[font=\tiny, below left] at (0,0) {$O$};
  % 抛物线 y^2=-4x，即 x=-t^2/4
  \draw[black, thick, domain=-2.0:2.0, smooth, samples=60]
    plot ({-(\x)^2/4}, \x);
  % 焦点
  \fill[black] (-1,0) circle [radius=2pt];
  \node[below left, font=\tiny] at (-1,0) {$F$};
  % 准线
  \draw[red, thin] (1,-2.1) -- (1,2.1);
  % 标注
  \node[font=\small, align=center] at (0,-2.8) {$y^2=-2px$\\开口向左};
\end{scope}

%% --- 子图3：开口上 x^2 = 2py (p=2, x^2=4y) ---
\begin{scope}[shift={(0,-5.5)}]
  \draw[->, thin, gray] (-2.3,0) -- (2.3,0) node[right, font=\tiny] {$x$};
  \draw[->, thin, gray] (0,-1.5) -- (0,3.2) node[above, font=\tiny] {$y$};
  \node[font=\tiny, below right] at (0,0) {$O$};
  % 抛物线 x^2=4y，即 y=t^2/4
  \draw[black, thick, domain=-2.0:2.0, smooth, samples=60]
    plot (\x, {(\x)^2/4});
  % 焦点
  \fill[black] (0,1) circle [radius=2pt];
  \node[right, font=\tiny] at (0.05,1) {$F$};
  % 准线
  \draw[red, thin] (-2.1,-1) -- (2.1,-1);
  % 标注
  \node[font=\small, align=center] at (0,-2.0) {$x^2=2py$\\开口向上};
\end{scope}

%% --- 子图4：开口下 x^2 = -2py (p=2, x^2=-4y) ---
\begin{scope}[shift={(5.5,-5.5)}]
  \draw[->, thin, gray] (-2.3,0) -- (2.3,0) node[right, font=\tiny] {$x$};
  \draw[->, thin, gray] (0,-3.2) -- (0,1.5) node[above, font=\tiny] {$y$};
  \node[font=\tiny, above right] at (0,0) {$O$};
  % 抛物线 x^2=-4y，即 y=-t^2/4
  \draw[black, thick, domain=-2.0:2.0, smooth, samples=60]
    plot (\x, {-(\x)^2/4});
  % 焦点
  \fill[black] (0,-1) circle [radius=2pt];
  \node[right, font=\tiny] at (0.05,-1) {$F$};
  % 准线
  \draw[red, thin] (-2.1,1) -- (2.1,1);
  % 标注
  \node[font=\small, align=center] at (0,-2.7) {$x^2=-2py$\\开口向下};
\end{scope}

% 总标题
\node[font=\normalsize, align=center] at (2.75, 2.9)
  {四种标准抛物线（$p>0$）};

\end{tikzpicture}
\end{document}
